\begin{table*}[t]
\centering
\small
\begin{tabular}{@{}l>{\raggedright\arraybackslash}p{4.1cm}>{\centering\arraybackslash}p{1.7cm}>{\centering\arraybackslash}p{1.9cm}>{\centering\arraybackslash}p{1.7cm}>{\centering\arraybackslash}p{2.2cm}@{}}
\toprule
Method & Run status & Feasible traces & Full-trace matches & Explicit repairs & Selected repair enforced \\
\midrule
OptABA-PC & \textbf{400/400 saved}: 368 completed, 32 timeout incumbents & 88/400 & 88/400 & \textbf{312/312} & \textbf{400/400} \\
ABA-PC & \textbf{400/400 saved and completed} & 88/400 & 88/400 & \textbf{312/312} & \textbf{400/400} \\
ASPCR-DAG & \textbf{250/250 saved; all optima proved} & 32/250 & 32/250 & -- & -- \\
MPC & \textbf{400/400 saved and completed} & 88/400 & 80/400 & -- & -- \\
FGS & \textbf{400/400 saved and completed} & -- & -- & -- & -- \\
\bottomrule
\end{tabular}
\caption{Evidence correspondence, run completion, and repair coverage. A trace is feasible if some DAG satisfies every recorded CI outcome; a full-trace output is the method's returned graph or class satisfying all of them. Thus MPC has 88 feasible inputs but only 80 matching outputs. A timeout incumbent is a valid saved repair whose optimality was not proved before the limit. An explicit repair names excluded facts, and selected-repair enforcement verifies every retained fact against the output. ASPCR-DAG uses its primary matched run set; FGS has no CI trace. A dash denotes an undefined object; bold marks complete audited coverage, not significance.}
\label{tab:contestability-overview}
\end{table*}
